\section{Tulokset}
Merkitään konsistenssikokeen kysymykset A-E ja herkkyyskokeen variaatiot A1-A4, B1-B4, C1-C4, D1-D4 ja E1-E4 listojen \ref{it:q_consistency} ja \ref{it:q_sensitivity} mukaisesti. Kokeiden tulokset on esitetty taulukoissa \ref{tb:consistency} ja \ref{tb:sensitivity_all}. Taulukot kuvaavat kunkin kysymyksen osalta, kuinka monta kertaa järjestelmän tuottama SQL-kysely oli korrekti, melkein korrekti, virheellinen tai johti tilanteeseen, jossa SQL-kyselyä ei tuotettu lainkaan.

Konsistenssikokeessa vain kysymys C tuotti sekä korrekteja että virheellisiä SQL-kyselyitä. Kysymys B tuotti kaksi uniikkia SQL-kyselyä, joista toinen käytti semanttisessa näkymässä määriteltyä metriikkaa \texttt{answer\_percentage} suoraan ja toinen laski vastaavan prosenttiluvun erillisellä laskukaavalla.

Kysymykset A, D ja E tuottivat konsistenssikokeessa vain yhden uniikin lopputuloksen kymmenessä suorituksessa. Vaikka kysymys E tuotti jokaisella suorituskerralla virheellisen vastauksen siinä mielessä, että SQL-kyselyä ei muodostettu, konsistenssin näkökulmasta järjestelmä oli johdonmukainen, koska tuloste oli joka kerta sama. Kysymys E osoittaa, että konsistenssi ei yksin riitä kuvaamaan järjestelmän luotettavuutta.

Tulokset viittaavat siihen, että tässä tutkimusasetelmassa havaitut ongelmat olivat usein toistettavia, mikä on hyödyllinen havainto sekä käyttäjän että semanttisen näkymän kehittäjän näkökulmasta. Kysymys C kuitenkin osoittaa, että järjestelmä pystyi tuottamaan samassa asetelmassa sekä korrekteja että virheellisiä vastauksia, joten toistettavuus ei ollut täydellistä.

Alla on esimerkki kysymykseen E tuotetusta ei kyselyä -vastauksesta:
\begin{quote}
  I'm sorry, but this question cannot be fully answered with the available data. While we can determine how many callbacks were answered versus not answered (using the c\_answered field filtered for callback media type), the semantic model does not contain information about the original inbound contact that triggered a callback. There is no field linking a callback conversation back to its originating inbound contact, so we cannot calculate the time between the original inbound contact and the callback attempt.
\end{quote}

Kysymykseen B järjestelmä tuotti kaksi uniikkia SQL-kyselyä. Näistä toinen käytti semanttisessa näkymässä ennalta määriteltyä metriikkaa \texttt{answer\_percentage} ja toinen muodosti vastaavan laskukaavan eksplisiittisesti SQL:ssä. Kysymykseen C järjestelmä tuotti kaksi uniikkia SQL-kyselyä, joista toinen oli korrekti ja toinen virheellinen. Korrekti kysely tuotettiin seitsemän kertaa ja virheellinen kolme kertaa. Virheellisessä kyselyssä järjestelmä sekoitti asiakkaan ja agentin roolit keskenään.

Konsistenssikokeen 50 suorituksesta 37 tuotti korrektin SQL-kyselyn, 10 johti tilanteeseen, jossa järjestelmä ei tuottanut lainkaan SQL-kyselyä, ja 3 tuotti virheellisen SQL-kyselyn. Tulokset osoittavat, että Cortex Analyst pystyy usein tuottamaan johdonmukaisia SQL-kyselyitä, mutta sen luotettavuus ei ollut kaikissa kysymystyypeissä riittävällä tasolla.

Herkkyystesteissä suurin variaatio havaittiin kysymyksen C eri muotoiluissa, joissa järjestelmä tuotti yhteensä kuusi uniikkia SQL-kyselyä. Kun Cortex Analystille syötettiin kysymyksen C eri variaatioita, järjestelmä tuotti väärän vastauksen 14 kertaa 50 suorituksesta. Virheellisissä kyselyissä toistuvin ongelma oli se, että järjestelmä ei rajannut hakua sarakkeen \texttt{direction} mukaisesti, vaan käytti saraketta \texttt{direction\_original} tai jätti suuntasuodatuksen kokonaan pois.

Kysymyksessä C ja sen variaatioissa puhutaan palvelutasosta, jolloin implisiittisenä taustaoletuksena on, että tarkastelu kohdistuu vain saapuviin kontakteihin. Havainto viittaa siihen, että Cortex Analyst ei tässä tutkimusasetelmassa aina huomioinut tällaisia implisiittisiä oletuksia johdonmukaisesti.

Kun kysymys E ja sen variaatiot annettiin järjestelmälle syötteeksi, järjestelmä tuotti 40 kertaa 50 suorituksesta tekstivastauksen, jossa se ilmoitti, että kysymykseen ei voida vastata. Variaatio E2 tuotti yhdeksän kertaa korrektin SQL-kyselyn ja kerran kyselyn, joka oli muuten oikea mutta käytti \texttt{distinct(identifier)}-ehdon sijaan ehtoa, jossa koko rivi vaadittiin uniikiksi. Tämä on huomattavasti heikompi ehto ja johti duplikaattivirheisiin.

Kysymyksen E variaatio E2 ei implikoi yhtä vahvasti, että kysymykseen vastaaminen vaatisi kahden rivin välistä vertailua. Tämä voi osaltaan selittää, miksi juuri tämä variaatio tuotti enemmän oikeita vastauksia kuin muut.

Kysymys E oli Cortex Analystille vaikein, mikä painottaa virhetyyppien jakaumaa ei kyselyä -luokan suuntaan.

Taulukosta \ref{tb:errors} nähdään, että muista virhetyypeistä yleisimpiä olivat suodatusvirheet ja duplikaattien hallintaan liittyvät virheet. Kaikki suodatusvirheet liittyivät sarakkeisiin \texttt{direction} ja \texttt{direction\_original}, mikä viittaa siihen, että järjestelmä voisi olla robustimpi, jos nämä kaksi saraketta erotettaisiin toisistaan semanttisesti selvemmin. Esimerkiksi \texttt{direction\_original}-sarakkeen korvaaminen \texttt{is\_callback}-sarakkeella saattaisi helpottaa mallia erottamaan sarakkeiden merkitykset toisistaan. Tätä ei kuitenkaan testattu tässä tutkimuksessa.

On kuitenkin huomioitava, että juuri nämä sarakkeet nousivat esiin tässä tutkimusasetelmassa. Toisenlaisessa tietokannassa tai toisessa sovellusympäristössä ongelmalliset kentät voisivat olla erilaisia. Havainto korostaa semanttisen näkymän merkitystä järjestelmän luotettavuuden kannalta. Semanttisen näkymän huolellinen suunnittelu ja iteratiivinen kehittäminen voivat pitkällä aikavälillä parantaa järjestelmän toimintaa merkittävästi.

Duplikaattien hallintavirheillä tarkoitetaan tässä virheitä, joissa \texttt{distinct}-avainsanaa on käytetty tarpeettomasti tai väärin. Tutkimusasetelmassa vain kysymystä A vastaavissa SQL-kyselyissä \texttt{distinct}-avainsanaa ei tarvita. Muiden kysymysten kohdalla duplikaattien hallinta tulisi yleensä kohdistaa \texttt{identifier}-sarakkeeseen. Duplikaattien hallintavirheitä esiintyi tutkimuksessa yhteensä 18 kertaa 250 suorituksen aikana.

Kysymyksessä A duplikaattien hallintaa ei odoteta lainkaan, mutta 10 suorituksessa järjestelmä tuotti SQL-kyselyn, joka sisälsi lausekkeen \texttt{distinct(identifier)}. Muut havaitut duplikaattien hallintavirheet liittyivät siihen, että \texttt{distinct}-avainsanaa sovellettiin n-tupleen, joka sisälsi identifier-sarakkeen lisäksi myös muita sarakkeita. Tällöin kysely vaatii näiden sarakkeiden yhdistelmän olevan uniikki, mikä on heikompi ehto kuin se, että \texttt{identifier}-sarake olisi jo yksin uniikki.

Taululiitos- tai sarakevalintavirheitä ei havaittu. Tämä viittaa siihen, että järjestelmä tunnisti kysymyksistä yleensä oikein, mitä attribuutteja käyttäjä halusi hakea, vaikka se ei aina onnistunut tuottamaan oikeaa SQL-kyselyä niiden perusteella.

\renewcommand{\arraystretch}{1.1}
\begin{table}[h]
  \centering
  \caption{Konsistenssikokeen tulokset.}
  \label{tb:consistency}
  \setlength{\tabcolsep}{4pt}
  {\small
    \begin{tabular}{lccccc}
      \toprule
      Kysymys & Oikein & \makecell{Melkein\\oikein} & Väärin & \makecell{Ei\\kyselyä} &
      \makecell{Erilaisten SQL-\\kyselyiden määrä} \\
      \midrule
      A & 10 & 0 & 0 & 0 & 1 \\
      B & 10 & 0 & 0 & 0 & 2 \\
      C & 7  & 0 & 3 & 0 & 2 \\
      D & 10 & 0 & 0 & 0  & 1 \\
      E & 0  & 0 & 0 & 10 & 1 \\
      \bottomrule
    \end{tabular}
  }
\end{table}

\begin{table}[h]
  \centering
  \caption{Herkkyystestien 1--5 tulokset. Sarakkeet ilmoittavat oikeiden, pienen virheen sisältävien, väärien ja ei kyselyä -vastausten lukumäärät sekä yhteenlaskettu uniikkien SQL-kyselyiden määrä. (\textit{Yht.}).}
  \label{tb:sensitivity_all}

  \centering
  \setlength{\tabcolsep}{4pt}
  \begin{tabular}{lccccc}
    \toprule
    Kysmys & Oikein & \makecell{Melkein\\oikein} & Väärin & \makecell{Ei\\kyselyä} & Yht. \\
    \midrule
    A  & 10 & 0  & 0  & 0 & 1 \\
    A1 & 10 & 0  & 0  & 0 & 1 \\
    A2 & 10 & 0  & 0  & 0 & 1 \\
    A3 & 10 & 0  & 0  & 0 & 1 \\
    A4 & 0  & 10 & 0  & 0 & 2 \\
    \bottomrule
  \end{tabular}

  \vspace{0.3em}
  \centering
  \setlength{\tabcolsep}{4pt}
  \begin{tabular}{lccccc}
    \toprule
    Kysmys & Oikein & \makecell{Melkein\\oikein} & Väärin & \makecell{Ei\\kyselyä} & Yht. \\
    \midrule
    B  & 10 & 0 & 0 & 0 & 2 \\
    B1 & 10 & 0 & 0 & 0 & 2 \\
    B2 & 10 & 0 & 0 & 0 & 2 \\
    B3 & 10 & 0 & 0 & 0 & 2 \\
    B4 & 10 & 0 & 0 & 0 & 2 \\
    \bottomrule
  \end{tabular}

  \vspace{0.3em}
  \centering
  \setlength{\tabcolsep}{4pt}
  \begin{tabular}{lccccc}
    \toprule
    Kysmys & Oikein & \makecell{Melkein\\oikein} & Väärin & \makecell{Ei\\kyselyä} & Yht. \\
    \midrule
    C  & 7 & 0  & 3  & 0 & 2 \\
    C1 & 8 & 2  & 0  & 0 & 3 \\
    C2 & 10 & 0  & 0 & 0 & 3 \\
    C3 & 0 & 0 & 10  & 0 & 5 \\
    C4 & 8 & 1  & 1  & 0 & 6 \\
    \bottomrule
  \end{tabular}

  \vspace{0.3em}
  \centering
  \setlength{\tabcolsep}{4pt}
  \begin{tabular}{lccccc}
    \toprule
    Kysmys & Oikein & \makecell{Melkein\\oikein} & Väärin & \makecell{Ei\\kyselyä} & Yht. \\
    \midrule
    D  & 10 & 0 & 0 & 0 & 1 \\
    D1 & 10 & 0 & 0 & 0 & 1 \\
    D2 & 10 & 0 & 0 & 0 & 1 \\
    D3 & 10 & 0 & 0 & 0 & 2 \\
    D4 & 7  & 0 & 3 & 0 & 3 \\
    \bottomrule
  \end{tabular}

  \vspace{0.3em}
  \centering
  \setlength{\tabcolsep}{4pt}
  \begin{tabular}{lccccc}
    \toprule
    Kysmys & Oikein & \makecell{Melkein\\oikein} & Väärin & \makecell{Ei\\kyselyä} & Yht. \\
    \midrule
    E  & 0 & 0 & 0 & 10 & 1 \\
    E1 & 0 & 0 & 0 & 10 & 1 \\
    E2 & 9 & 1 & 0 & 0 & 3 \\
    E3 & 0 & 0 & 0 & 10 & 3 \\
    E4 & 0 & 0 & 0 & 10 & 3 \\
    \bottomrule
  \end{tabular}
\end{table}

\begin{table}[h]
  \centering
  \caption{Esiintyneiden virhetyyppien määrät kaikissa kokeissa. Rivit kuvaavat, kuinka monta kertaa kukin virhetyyppi esiintyi kussakin kysymyksessä tai sen variaatioissa. Jos yhdessä SQL-kyselyssä esiintyy useampi virhetyyppi, kysely lasketaan jokaiseen virhetyyppiin erikseen.}
  \label{tb:errors}
  \setlength{\tabcolsep}{4pt}
  \begin{tabular}{lcccccc}
    \toprule
    Kysymys & A & B & C & D & E & Yhteensä\\
    \midrule
    Aggregaatio & 0 & 0 & 2  & 0 & 0  & 2 \\
    Duplikaattien hallinta & 10 & 0 & 7 & 0 & 1 & 18 \\
    Ei kyselyä & 0 & 0 & 0 & 0 & 40 & 40 \\
    Sarakevalinta & 0 & 0 & 0 & 0 & 0 & 0 \\
    Suodatus & 0 & 0 & 11  & 3 & 0  & 14 \\
    Taululiitos & 0 & 0 & 0 & 0 & 0 & 0 \\
    \bottomrule
  \end{tabular}
\end{table}
